% !TeX program = pdflatex
\documentclass[11pt,a5paper]{article}
\usepackage[french]{babel}
\usepackage[T1]{fontenc}
\usepackage[utf8]{inputenc}
\usepackage{ebgaramond}
\usepackage[a5paper,top=16mm,bottom=18mm,inner=16mm,outer=16mm,headheight=10pt,headsep=6mm,footskip=10mm]{geometry}
\usepackage{xcolor,graphicx,tikz,microtype,fancyhdr,ragged2e,needspace,pdfpages}
\definecolor{bachred}{HTML}{8A2822}
\definecolor{ink}{HTML}{242220}
\definecolor{muted}{HTML}{6F6660}
\color{ink}
\setlength{\parindent}{0pt}
\setlength{\parskip}{0.43em}
\linespread{1.03}
\emergencystretch=2em
\pagestyle{fancy}
\fancyhf{}
\fancyhead[C]{\fontsize{8}{9}\selectfont\color{muted}\scshape Bach\,\,\textbullet\,\, Suites pour violoncelle seul}
\fancyfoot[C]{\color{bachred}\fontsize{9}{10}\selectfont\thepage}
\renewcommand{\headrulewidth}{0pt}
\renewcommand{\footrulewidth}{0pt}
\newcommand{\divider}{\par\vspace{3mm}\begin{center}\color{bachred}\rule{36mm}{0.45pt}\end{center}\vspace{1.3mm}}
\newcommand{\smallcapshead}[1]{\noindent{\color{bachred}\fontsize{13.7}{15.5}\selectfont\bfseries #1}\par\vspace{1.7mm}}
\newcommand{\pagetitle}[2]{\begin{center}{\color{bachred}\fontsize{9}{11}\selectfont\textls[115]{\MakeUppercase{#1}}}\\[1.7mm]{\fontsize{19}{20.8}\selectfont #2}\\[1mm]{\color{bachred}\rule{40mm}{0.45pt}}\end{center}\vspace{3.1mm}}
\newcommand{\suite}[3]{\needspace{36mm}{\color{bachred}\fontsize{15}{17}\selectfont Suite n\textsuperscript{o} #1}\hfill{\color{muted}\fontsize{10}{12}\itshape #3}\par{\fontsize{10}{12}\itshape #2}\par\vspace{1.3mm}}
\newcommand{\suiteitems}[1]{{\fontsize{10.5}{13.1}\selectfont #1}\par\vspace{4.1mm}}
\begin{document}
\includepdf[pages=1]{couverture_ordree.pdf}
\setcounter{page}{2}
\pagetitle{Récital}{Programme musical}
\begin{center}{\fontsize{11}{13}\selectfont Johann Sebastian Bach}\par{\color{muted}\fontsize{9.5}{11}\selectfont (1685--1750)}\end{center}\vspace{2mm}
\suite{3}{BWV 1009}{ut majeur}
\suiteitems{I. Prélude\quad\textbullet\quad II. Allemande\par III. Courante\quad\textbullet\quad IV. Sarabande\par V. Bourrées I et II\quad\textbullet\quad VI. Gigue}
\suite{2}{BWV 1008}{ré mineur}
\suiteitems{I. Prélude\quad\textbullet\quad II. Allemande\par III. Courante\quad\textbullet\quad IV. Sarabande\par V. Menuets I et II\quad\textbullet\quad VI. Gigue}
\suite{1}{BWV 1007}{sol majeur}
\suiteitems{I. Prélude\quad\textbullet\quad II. Allemande\par III. Courante\quad\textbullet\quad IV. Sarabande\par V. Menuets I et II\quad\textbullet\quad VI. Gigue}
\vfill\begin{center}{\color{muted}\fontsize{9.5}{11}\itshape Lucien Hennecart \quad\textbullet\quad violoncelle}\end{center}
\newpage
\pagetitle{Pour accompagner l'écoute}{Les Suites de Bach}
\smallcapshead{Une musique pour un seul instrument}
\begin{justify}
Les six \emph{Suites pour violoncelle seul} de Jean-Sébastien Bach ont vraisemblablement été composées pendant ses années à Köthen, autour de 1720. Sans accompagnement, le violoncelle y suggère tour à tour une mélodie, une basse, des accords et plusieurs voix. L'instrument semble parfois jouer à lui seul toute une conversation musicale.

Chacune des Suites s'ouvre par un prélude et se poursuit par cinq mouvements de danse : allemande, courante, sarabande, deux danses de même nom, puis gigue. Leurs caractères passent de la gravité à la légèreté, de l'intimité à l'énergie.
\end{justify}
\divider
\smallcapshead{Trois Suites, trois caractères}
\begin{justify}
La \emph{troisième Suite}, en ut majeur, ouvre ce récital avec sa sonorité ample, ses accords affirmés et l'élan de ses Bourrées. Plus retenue, la \emph{deuxième Suite}, en ré mineur, fait entendre une écriture intérieure, parfois presque dépouillée, particulièrement dans sa Sarabande.

La \emph{première Suite}, en sol majeur, clôt le programme. Son Prélude, devenu l'une des pages les plus familières du répertoire pour violoncelle, est suivi de danses au caractère clair et souple.
\end{justify}
\newpage
\pagetitle{L'interprète}{Lucien Hennecart}
\noindent
\begin{minipage}[t]{0.45\textwidth}
\vspace{0pt}\includegraphics[width=\linewidth,height=78mm,keepaspectratio]{photo-lucien.jpg}
\end{minipage}\hfill
\begin{minipage}[t]{0.51\textwidth}
\vspace{0pt}\fontsize{9.4}{11.6}\selectfont
Lucien Hennecart est violoncelliste et chercheur en mathématiques. Depuis décembre 2024, il est chargé de recherche au CNRS, au LAMFA, à l'Université de Picardie Jules Verne. Après une thèse à l'Université Paris-Saclay, il a poursuivi ses recherches à Édimbourg, où il a aussi joué au sein de l'orchestre universitaire.
\end{minipage}
\vspace{3mm}
\smallcapshead{Parcours musical}
\begin{justify}
\fontsize{9.9}{12.2}\selectfont
Il a étudié le violoncelle au conservatoire de Saint-Étienne auprès de Marianne Gaiffe et Pascal Jemain, ainsi qu'à Rennes auprès d'Alain Brunier. Il a participé à des masterclasses avec Emmanuelle Bertrand, Jean Deplace, Xavier Phillips et Claudio Pasceri, et suivi des stages avec Xavier Gagnepain.

À Amiens, il joue au sein de l'Orchestre universitaire de Picardie et de son Ensemble baroque.
\end{justify}
\vspace{1.6mm}
{\color{bachred}\rule{\linewidth}{0.45pt}}\par
\begin{justify}\fontsize{9.3}{11.6}\selectfont
\emph{Son instrument.} Lucien Hennecart joue un violoncelle réalisé en 2007 par Geoffroy Mercier, luthier à Rontalon, dans le Rhône, dont le travail s'appuie notamment sur l'étude des tracés géométriques des instruments du quatuor.
\end{justify}
\vfill\begin{center}{\color{bachred}\itshape Merci de votre présence et belle écoute !}\end{center}
\end{document}
